% aspectratio=169 must stay on \documentclass: it is a Beamer class option
% and cannot be set by the theme.
\documentclass[aspectratio=169]{beamer}
\usetheme{ui1beamer}

% ==============================================================
% FILL IN YOUR DETAILS BELOW
% ==============================================================
\title{Valoración de inversiones}
\subtitle{Métodos VAN y TIR}
\asignatura{Matemáticas Financieras}
\author{Israel Herraiz Tabernero}
\date{9 de agosto de 2026}
% ==============================================================

\begin{document}

\begin{frame}[plain]
  \titlepage
\end{frame}

\section{Introducción}

\begin{frame}{Objetivos de la sesión}
  \begin{itemize}
    \item Repasar el concepto de flujo de caja.
    \item Calcular el valor actual neto (VAN) de un proyecto.
    \item Interpretar la tasa interna de retorno (TIR).
    \item Comparar proyectos con horizontes temporales distintos.
  \end{itemize}

  \begin{block}{Idea principal}
    Un euro de hoy vale más que un euro de mañana: toda la valoración de
    inversiones se deriva de esta observación.
  \end{block}
\end{frame}

\begin{frame}{Estructura de la unidad}
  \begin{enumerate}
    \item Flujos de caja y descuento.
    \item Criterios de decisión.
    \item Casos prácticos.
  \end{enumerate}
\end{frame}

\section{Criterios de decisión}

\begin{frame}{Valor actual neto}
  El VAN descuenta cada flujo de caja \(F_t\) a la tasa \(k\):

  \begin{equation*}
    \mathrm{VAN} = -F_0 + \sum_{t=1}^{n} \frac{F_t}{(1+k)^t}
  \end{equation*}

  \begin{itemize}
    \item \(\mathrm{VAN} > 0\): el proyecto crea valor.
    \item \(\mathrm{VAN} < 0\): el proyecto destruye valor.
  \end{itemize}
\end{frame}

\begin{frame}{Tasa interna de retorno}
  \begin{columns}[T]
    \begin{column}{0.5\textwidth}
      La TIR es la tasa que anula el VAN:
      \begin{equation*}
        -F_0 + \sum_{t=1}^{n} \frac{F_t}{(1+\mathrm{TIR})^t} = 0
      \end{equation*}
    \end{column}
    \begin{column}{0.5\textwidth}
      \begin{exampleblock}{Regla de decisión}
        Aceptar el proyecto si la TIR supera el coste de capital.
      \end{exampleblock}
    \end{column}
  \end{columns}

  \begin{alertblock}{Cuidado}
    Con flujos de signo alternante puede haber varias TIR, o ninguna.
  \end{alertblock}
\end{frame}

\section{Caso práctico}

\begin{frame}{Comparación de dos proyectos}
  \begin{center}
    \begin{tabular}{lrr}
      \textbf{Año} & \textbf{Proyecto A} & \textbf{Proyecto B} \\
      \hline
      0 & $-10\,000$ & $-10\,000$ \\
      1 & $6\,000$    & $2\,000$    \\
      2 & $5\,000$    & $4\,000$    \\
      3 & $1\,000$    & $8\,000$    \\
      \hline
      VAN (k = 8\,\%) & $1\,005$ & $1\,318$ \\
    \end{tabular}
  \end{center}

  A igualdad de inversión inicial, B crea más valor pese a recuperar la
  inversión más tarde.
\end{frame}

\begin{frame}{Resumen}
  \begin{itemize}
    \item El VAN mide creación de valor en euros de hoy.
    \item La TIR mide rentabilidad relativa, pero es frágil.
    \item Ante conflicto entre criterios, prevalece el VAN.
  \end{itemize}

  \begin{block}{Para la próxima sesión}
    Resolver los ejercicios 1 a 5 de la unidad didáctica.
  \end{block}
\end{frame}

\end{document}
